\SourceLeaf{231}%
\bigskip\Needspace{9\baselineskip}
\pdfbookmark[1]{Pallas Olbersiana. 1802 Oktober}{pallas180210}
\begin{center}
PALLAS OLBERSIANA.\par
\medskip\hrule\medskip
Monatliche Correspondenz. Band VI. S.\ 390—396. 1802 Oktober.\par
\medskip\hrule
\end{center}\medskip
—\quad —\quad —

Dr.\ \textsc{Gauß}, der die Elemente der Pallasbahn zum dritten Mal verbessert hatte, hat nun eine \textit{vierte} Korrektur gewagt. Diese neuen Elemente sind aus unseren ersten Beobachtungen auf dem \textit{Seeberg}, aus Dr.\ \textsc{Maskelynes} letzten Meridianbeobachtungen in \textit{Greenwich} und aus Dr.\ \textsc{Olbers}’ Beobachtung vom 8. Juli abgeleitet. Dr.\ \textsc{Gauß} möchte sie aber noch nicht als zuverlässiger ausgeben als die Elemente III, da er von den späteren Mailänder Beobachtungen eine viel genauere Verbesserung erwartet. Die Elemente lauten indessen wie folgt. [Dieselben wie oben im Brief an \textsc{Bode} vom 10. September, S.\ 223, ebenfalls für die Epoche 31. März 1802, mittags auf dem Seeberg, mit dem Zusatz „für die Epoche, siderisch ruhend“ beim Aphel und beim Knoten.]

Mit diesen Elementen hat Dr.\ \textsc{Gauß} die folgende Ephemeride für die Bewegung des Planeten im kommenden Jahr vorausberechnet. Es ist sinnvoll, sie frühzeitig bekannt zu machen. So kann man in der Zwischenzeit in dem Gebiet, das \textit{Pallas} im Jahr 1803 durchwandern wird, zweifelhafte Sterne genauer bestimmen, sich in diesem Himmelsgebiet vorab orientieren und sich damit vertraut machen. Das wird uns das Auffinden erleichtern. In der letzten Spalte der Ephemeride steht die Lichtstärke des Planeten. Als Einheit dient die Lichtstärke, die er am 8. Juli dieses Jahres hatte. Nach demselben Maßstab betrug sie:
\begin{center}
\begin{tabular}{r l l r}
am 4. April&1802&$=$&$4.23$\\
4. Mai&—&&$2.75$\\
16. Juli&—&&$0.90$\\
26. Juli&—&&$0.80$\\
7. Aug.&—&&$0.70$\\
24. Aug.&—&&$0.60$
\end{tabular}
\end{center}
Die Größe der Phase, die Höhe über dem Horizont und der zeitliche Abstand von der Dämmerung sind hierbei unberücksichtigt geblieben. Der erste Umstand ist ganz unerheblich, die beiden anderen hängen von den Beobachtungsorten und den leistungsstärkeren Fernrohren ab. Die größte Lichtstärke, die \textit{Pallas} im kommenden Jahr erreicht, wird $=0.656$ sein. Dieselbe Lichtstärke hatte sie in diesem Jahr 1802 am 14. August. \textsc{Oriani} \textit{beobachtete} den Planeten noch am 8. August und \textit{sah} ihn mit Mühe am 17. und 18. desselben Monats, wozu die Dämmerung und der niedrige Stand des Planeten wahrscheinlich viel beitrugen. So bleibt uns noch immer die angenehme Hoffnung, dass wir \textit{Pallas} im kommenden Jahr zu Gesicht bekommen und ihre Opposition zur Sonne beobachten können. Wahrscheinlich wird der Planet schon vor diesem Zeitpunkt (30. Juni 1803) gefunden, denn mit dem Zusammentreffen nimmt seine Sichtbarkeit oder Lichtstärke schon wieder ab. Im Meridian werden wir wohl, zumindest in unseren Breiten, vor der zweiten Maihälfte oder vor Anfang Juni wenig Hoffnung haben, diesen Planeten zu beobachten. [Die Ephemeride reicht vom 4. Februar bis zum 28. Juni, Mitternacht auf dem Seeberg 1803. Sie gibt für jeden dritten Tag RA und Deklination auf Minuten genau, den Abstand von der Erde auf zwei und die Lichtstärke auf drei Dezimalstellen an.]

Selbst wenn es Dr.\ \textsc{Gauß} gelingt, die obigen Elemente anhand der italienischen Beobachtung weiter zu verbessern, wie es höchst wahrscheinlich ist, wird es nicht nötig sein, die mühsame
\SourceLeaf{232}%
Berechnung einer neuen Ephemeride ganz auf dieser Grundlage zu wiederholen. Es wird dann genügen, \textit{einige} Positionen des Planeten nachzurechnen, denn der Unterschied kann nicht sehr beträchtlich sein und nur in langsamen Stufen gleichmäßig anwachsen. Dr.\ \textsc{Gauß} hat bereits versucht, einige Positionen der \textit{Pallas} nach seinen Elementen III zu berechnen. Der Unterschied ist kaum so beträchtlich, dass er das Auffinden des Planeten behindern dürfte, wie aus dem folgenden Vergleich zu sehen ist. [Diesem sind hier die Ergebnisse für die Elemente V aus der Handschrift beigefügt.]
\begin{center}
\begingroup\small\setlength{\tabcolsep}{2.5pt}
\begin{tabular}{l|r|r|r|r|r|r|r|r}
&\multicolumn{2}{c|}{RA der $\PallasSymbol$}&&\multicolumn{2}{c|}{Dekl. der $\PallasSymbol$}&&RA&Dekl.\\
1803&III. Elem.&IV. Elem.&Diff.&III. Elem.&IV. Elem.&Diff.&V. Elem.&V. Elem.\\\hline
Febr.\ 4&$267^0\ 35'$&$267^0\ 48'$&$13'$&$5^0\ 41'$&$5^0\ 38'$&$3'$&$267^0\ 45'$&$5^0\ 39'$\\
Juni 28&$275\ 45$&$276\ 17$&$32$&$23\ 12$&$23\ 10$&$2$&$276\ 9$&$23\ 11$
\end{tabular}\endgroup
\end{center}
Die Opposition dieses Planeten fällt nach den Elementen III auf den Vormittag des 30. Juni bürgerlicher Zeit, nach den Elementen IV auf den Nachmittag desselben Tages. Der Unterschied in der Deklination ist ganz unbedeutend. Nach V. Juni 30. $2^{\mathrm h}\ 53^{\mathrm m}$ In $277^0\ 45'$ L.
\begin{flushright}[\textsc{von Zach.}]\end{flushright}

